\begin{textsample}{POS Dim 4 – human – Score 5.00 – mg/mg\_ef\_linguagens\_lp\_6\_p1.txt}  \label{ex:f4_pos_008}
Esse \textbf{eixo}, como já apresentado anteriormente neste documento e reforçado pela BNCC quando diz que esse \textbf{eixo} envolve os procedimentos e estratégias (meta)cognitivas de análise e avaliação consciente, durante os processos de \textbf{leitura} e de produção de \textbf{textos} ( \textbf{orais}, escritos e multissemióticos ), das materialidades dos \textbf{textos}, responsável por seus efeitos de sentido, seja no que se refere às formas de composição dos \textbf{textos}, determinadas pelos gêneros ( \textbf{orais}, escritos e multissemióticos ) e pela situação de produção, seja no que se refere aos estilos adotados nos \textbf{textos}, com forte impacto nos efeitos de sentido.
A análise linguística/semiótica tem como finalidade propor ao estudante a reflexão sobre as diferentes possibilidades e recursos da língua na produção de sentido e sua utilização nas práticas de interação social.

% matched lemmas: eixo, leitura, oral, texto
\end{textsample}
